\section{\texttt{aes.py}: function-by-function implementation}
Read conversion first, then field arithmetic, transformations, key expansion, encryption and inverse operations. Matrix transformations mutate their argument and return None; block functions return immutable bytes. The operation wrapper is the only generic instrumentation seam.
\subsection{\texttt{validate\_\allowbreak{}block}}
\begin{lightcode}
def validate_block(data, label="Block"):
\end{lightcode}
\noindent\textbf{Source:} complete \texttt{aes.py} listing in Part 6, lines 101--104.
\begin{enumerate}
\item Accept a byte sequence and a human-readable label. Check its length against BLOCK\_SIZE, which is sixteen.
\item Raise ValueError on a wrong length; otherwise return None. This validates size, not message padding or a general encryption mode.
\end{enumerate}
\subsection{\texttt{bytes\_\allowbreak{}to\_\allowbreak{}state}}
\begin{lightcode}
def bytes_to_state(data):
\end{lightcode}
\noindent\textbf{Source:} complete \texttt{aes.py} listing in Part 6, lines 107--116.
\begin{enumerate}
\item Validate the sixteen input elements and allocate four independent rows.
\item Traverse columns outside rows and assign data[4*column+row] to state[row][column]. Return a fresh mutable matrix, so later transformations do not mutate the caller's bytes.
\end{enumerate}
\subsection{\texttt{state\_\allowbreak{}to\_\allowbreak{}bytes}}
\begin{lightcode}
def state_to_bytes(state):
\end{lightcode}
\noindent\textbf{Source:} complete \texttt{aes.py} listing in Part 6, lines 119--121.
\begin{enumerate}
\item Traverse the matrix column by column in exactly the inverse conversion order.
\item Construct and return immutable bytes. For example, matrix entry state[1][2] becomes output byte nine, not byte six.
\end{enumerate}
\subsection{\texttt{xtime}}
\begin{lightcode}
def xtime(a):
\end{lightcode}
\noindent\textbf{Source:} complete \texttt{aes.py} listing in Part 6, lines 129--137.
\begin{enumerate}
\item Shift the byte left once to multiply its polynomial by x. Test the original high bit to detect the x\textasciicircum{}8 overflow.
\item If that bit was set, XOR 0x1B into the shifted result, then mask to eight bits. Thus xtime(0x80) is 0x1B, not the ordinary integer product 0x100.
\end{enumerate}
\subsection{\texttt{gf\_\allowbreak{}mul}}
\begin{lightcode}
def gf_mul(a, b):
\end{lightcode}
\noindent\textbf{Source:} complete \texttt{aes.py} listing in Part 6, lines 140--152.
\begin{enumerate}
\item Initialize an XOR accumulator and inspect eight bits of multiplier b. Whenever the low bit is one, XOR the current multiplicand a into the result.
\item Advance a with xtime and shift b right after each iteration. Return one field byte; multiplication by 02 and 03 in MixColumns therefore uses field reduction rather than ordinary multiplication.
\end{enumerate}
\subsection{\texttt{sub\_\allowbreak{}bytes}}
\begin{lightcode}
def sub_bytes(state):
\end{lightcode}
\noindent\textbf{Source:} complete \texttt{aes.py} listing in Part 6, lines 160--168.
\begin{enumerate}
\item Visit every state coordinate and use its current byte as an index into SBOX.
\item Overwrite that coordinate with the substituted byte. Return None: this operation changes values in the supplied matrix in place, without changing their positions.
\end{enumerate}
\subsection{\texttt{shift\_\allowbreak{}rows}}
\begin{lightcode}
def shift_rows(state):
\end{lightcode}
\noindent\textbf{Source:} complete \texttt{aes.py} listing in Part 6, lines 171--174.
\begin{enumerate}
\item For row indices one through three, split the row after its row index.
\item Concatenate the suffix and prefix to rotate that row left by its index. Row zero stays fixed, and the function updates the matrix in place without returning a new state.
\end{enumerate}
\subsection{\texttt{mix\_\allowbreak{}column}}
\begin{lightcode}
def mix_column(column):
\end{lightcode}
\noindent\textbf{Source:} complete \texttt{aes.py} listing in Part 6, lines 177--189.
\begin{enumerate}
\item Accept four original column bytes and initialize an empty output list. For each row of MIX\_MATRIX, XOR the four coefficient/byte products computed by gf\_mul.
\item Append each resulting byte and return the new four-byte list. Keeping the original column separate prevents an earlier output assignment from contaminating a later equation.
\end{enumerate}
\subsection{\texttt{mix\_\allowbreak{}columns}}
\begin{lightcode}
def mix_columns(state):
\end{lightcode}
\noindent\textbf{Source:} complete \texttt{aes.py} listing in Part 6, lines 192--197.
\begin{enumerate}
\item Extract one column into an independent four-element list and call mix\_column.
\item Write its four returned bytes back to the same column and repeat for all four columns. The state is updated in place; no printing or fault rule is introduced.
\end{enumerate}
\subsection{\texttt{add\_\allowbreak{}round\_\allowbreak{}key}}
\begin{lightcode}
def add_round_key(state, round_key):
\end{lightcode}
\noindent\textbf{Source:} complete \texttt{aes.py} listing in Part 6, lines 200--208.
\begin{enumerate}
\item Read the state and round-key matrices at matching row/column coordinates.
\item XOR each key byte into the state byte in place. Applying the same key twice undoes the operation, and applying it to both members of a differential pair preserves their XOR difference.
\end{enumerate}
\subsection{\texttt{rot\_\allowbreak{}word}}
\begin{lightcode}
def rot_word(word):
\end{lightcode}
\noindent\textbf{Source:} complete \texttt{aes.py} listing in Part 6, lines 216--218.
\begin{enumerate}
\item Take a four-byte key-schedule word and slice off its first byte.
\item Append that byte after the other three and return a new list. This rotates key-word bytes, not bits, and leaves the argument unchanged.
\end{enumerate}
\subsection{\texttt{sub\_\allowbreak{}word}}
\begin{lightcode}
def sub_word(word):
\end{lightcode}
\noindent\textbf{Source:} complete \texttt{aes.py} listing in Part 6, lines 221--223.
\begin{enumerate}
\item Take each byte of a key-schedule word as an SBOX index.
\item Return a new four-byte list of substituted values. This uses the same forward AES S-box as encryption, not INV\_SBOX.
\end{enumerate}
\subsection{\texttt{key\_\allowbreak{}expansion}}
\begin{lightcode}
def key_expansion(key):
\end{lightcode}
\noindent\textbf{Source:} complete \texttt{aes.py} listing in Part 6, lines 226--259.
\begin{enumerate}
\item Validate the master key and divide it into four four-byte words. Generate words four through forty-three by XORing the word four places back with a transformed copy of the previous word.
\item On every fourth word, first apply rot\_word, sub\_word and RCON to the first byte. Group each four consecutive words into a fresh state matrix and return eleven matrices K0 through K10.
\end{enumerate}
\subsection{\texttt{\_\allowbreak{}operation}}
\begin{lightcode}
def _operation(state, number, name, transform, observer):
\end{lightcode}
\noindent\textbf{Source:} complete \texttt{aes.py} listing in Part 6, lines 267--277.
\begin{enumerate}
\item Accept a mutable state, round number, operation name, transformation callable and optional observer. If present, notify the observer with a Before event and the live state.
\item Run transform(state), then issue the corresponding After event. The observer can modify the live matrix; the cipher itself defines no coordinate or XOR mask for an attack.
\end{enumerate}
\subsection{\texttt{aes\_\allowbreak{}encrypt}}
\begin{lightcode}
def aes_encrypt(plaintext, key, observer=None):
\end{lightcode}
\noindent\textbf{Source:} complete \texttt{aes.py} listing in Part 6, lines 280--291.
\begin{enumerate}
\item Convert the plaintext to a matrix, expand the key and apply the initial AddRoundKey with K0 through \_operation.
\item Run rounds one through ten in SubBytes, ShiftRows, optional MixColumns, AddRoundKey order. Skip MixColumns in round ten, serialize the final matrix and return ciphertext bytes without printing.
\end{enumerate}
\subsection{\texttt{inv\_\allowbreak{}sub\_\allowbreak{}bytes}}
\begin{lightcode}
def inv_sub_bytes(state):
\end{lightcode}
\noindent\textbf{Source:} complete \texttt{aes.py} listing in Part 6, lines 294--298.
\begin{enumerate}
\item Read each current state byte and index INV\_SBOX.
\item Replace the byte in place. The inverse table undoes the forward substitution independently at each coordinate.
\end{enumerate}
\subsection{\texttt{inv\_\allowbreak{}shift\_\allowbreak{}rows}}
\begin{lightcode}
def inv_shift_rows(state):
\end{lightcode}
\noindent\textbf{Source:} complete \texttt{aes.py} listing in Part 6, lines 301--304.
\begin{enumerate}
\item Visit rows one through three and split each row before its final row-index bytes.
\item Move those final bytes to the front. These right rotations undo shift\_rows while leaving row zero unchanged.
\end{enumerate}
\subsection{\texttt{inv\_\allowbreak{}mix\_\allowbreak{}columns}}
\begin{lightcode}
def inv_mix_columns(state):
\end{lightcode}
\noindent\textbf{Source:} complete \texttt{aes.py} listing in Part 6, lines 307--315.
\begin{enumerate}
\item Copy all four original bytes of each column before any assignment.
\item For every output row, XOR products with the corresponding INV\_MIX\_MATRIX coefficients, using gf\_mul, and store the result. The copy ensures every inverse output uses the same original column.
\end{enumerate}
\subsection{\texttt{aes\_\allowbreak{}decrypt}}
\begin{lightcode}
def aes_decrypt(ciphertext, key, observer=None):
\end{lightcode}
\noindent\textbf{Source:} complete \texttt{aes.py} listing in Part 6, lines 318--329.
\begin{enumerate}
\item Convert ciphertext to a state, expand the supplied key and remove K10 with AddRoundKey.
\item Iterate key indices nine down to zero, applying InvShiftRows, InvSubBytes and AddRoundKey, followed by InvMixColumns except at index zero. Serialize and return the recovered plaintext; this is ordinary decryption, not a key-recovery attack.
\end{enumerate}
\subsection{\texttt{main}}
\begin{lightcode}
def main():
\end{lightcode}
\noindent\textbf{Source:} complete \texttt{aes.py} listing in Part 6, lines 332--344.
\begin{enumerate}
\item Decode configured KEY and PLAINTEXT hexadecimal strings, encrypt once and decrypt the result.
\item Print the key, plaintext, ciphertext and recovered block; raise RuntimeError if the recovered input differs. The main guard calls this routine only when aes.py is executed directly.
\end{enumerate}

\subsection{Why the AddRoundKey lambdas are needed}
\texttt{\_\allowbreak{}operation} expects a transformation taking only the current state. An expression such as \texttt{lambda s: add\_\allowbreak{}round\_\allowbreak{}key(s, keys[number])} adapts the two-argument key-XOR primitive to that interface. It is invoked immediately during the current loop iteration; it is not stored for later execution. Consequently it uses the intended round key rather than a later loop index.
